\begin{theorem}
    (Пирсона) Если $\displaystyle H_{0}$ верна, то $\displaystyle \hat{\chi }_{n}^{2}\xrightarrow[n\rightarrow \infty ]{d} \chi _{m-1}^{2}$.
\end{theorem}
\begin{proof}
    Положим $\displaystyle Y_{j} =\begin{pmatrix}
    I_{X_{j} =a_{1}}\\
    \dotsc \\
    I_{X_{j} =a_{m}}
    \end{pmatrix} ,\ j=\overline{1,n}$ -- независимые одинаково распределенные случайные векторы, $\displaystyle EY_{j} =\begin{pmatrix}
    p_{1}\\
    \dotsc \\
    p_{m}
    \end{pmatrix} =\begin{pmatrix}
    p_{1}^{0}\\
    \dotsc \\
    p_{m}^{0}
    \end{pmatrix}$. Найдем матрицу ковариаций вектора $\displaystyle Y_{j}$:
    \begin{equation*}
        cov( I_{X_{j} =a_{i}} ,\ I_{X_{j} =a_{k}}) =EI_{X_{j} =a_{i} ,\ X_{j} =a_{k}} -EI_{X_{j} =a_{i}} \cdotp EI_{X_{j} =a_{k}} =\begin{cases}
        p_{i}^{0} -p_{i}^{0} p_{k}^{0} & ,\ i=k\\
        -p_{i}^{0} p_{k}^{0} & ,\ i\neq k
        \end{cases} .
    \end{equation*}
    Положим $\displaystyle B:=\begin{pmatrix}
    p_{1}^{0} & 0 & \dotsc  & 0\\
    0 & p_{2}^{0} & \dotsc  & 0\\
    \dotsc  & \dotsc  & \dotsc  & \dotsc \\
    0 & 0 & \dotsc  & p_{m}^{0}
    \end{pmatrix} \Rightarrow DY_{j} =B-\overline{p}_{0}\overline{p}_{0}^{T}$. По многомерной ЦПТ
    
    
    \begin{equation*}
        \dfrac{Y_{1} +\ \dotsc \ +Y_{n}}{\sqrt{n}} -\sqrt{n}\overline{p}_{0}\xrightarrow{d}\mathcal{N}\left( 0,\ B-\overline{p}_{0}\overline{p}_{0}^{T}\right) .
    \end{equation*}
    Заметим, что $\displaystyle Y_{1} +\ \dotsc \ +Y_{n} =\begin{pmatrix}
    \nu _{1}\\
    \dotsc \\
    \nu _{m}
    \end{pmatrix} =:\overline{\nu }$, и матрица $\displaystyle B$ симметрична и положительно определена в силу положительности $\displaystyle p_{j}^{0} \ \forall j\in \{1,\ \dotsc ,\ m\}$. Пусть $\displaystyle \xi _{n} =\left(\sqrt{B}\right)^{-1}\sqrt{n}\left(\dfrac{\overline{\nu }}{n} -\overline{p}_{0}\right)$. По теореме о наследовании сходимостей
    
    
    \begin{equation*}
        \xi _{n}\xrightarrow{d}\mathcal{N}\left( 0,\ \left(\sqrt{B}\right)^{-1}\left( B-\overline{p}_{0}\overline{p}_{0}^{T}\right)\left(\sqrt{B}\right)^{-1}\right) =\mathcal{N}\left( 0,\ I_{m} -zz^{T}\right) ,
    \end{equation*}
    где $\displaystyle z=\left(\sqrt{B}\right)^{-1}\overline{p}_{0} =\begin{pmatrix}
    \sqrt{p_{1}^{0}}\\
    \dotsc \\
    \sqrt{p_{m}^{0}}
    \end{pmatrix}$. Заметим, что $\displaystyle \Vert z\Vert ^{2} =\sum _{i=1}^{m} p_{i}^{0} =1\Rightarrow \Vert z\Vert =1$. Рассмотрим ортогональную матрицу $\displaystyle V\in M_{m\times m}$ такую, что ее первая строка -- это вектор $\displaystyle z$. Тогда по теореме о наследовании сходимости
    \begin{equation*}
        V\xi _{n}\xrightarrow{d} V\cdotp \mathcal{N}\left( 0,\ I_{m} -zz^{T}\right) =\mathcal{N}\left( 0,\ V\left( I_{m} -zz^{T}\right) V^{T}\right) .
    \end{equation*}
    Рассмотрим вектор $\displaystyle Vz$. В силу ортогональности матрицы $\displaystyle V$ выполнено
    \begin{equation*}
        Vz=\begin{pmatrix}
        \sqrt{p_{1}^{0}} & \sqrt{p_{2}^{0}} & \dotsc  & \sqrt{p_{m}^{0}}\\
        v_{21} & v_{22} & \dotsc  & v_{2m}\\
        \dotsc  & \dotsc  & \dotsc  & \dotsc \\
        v_{m1} & v_{m2} & \dotsc  & v_{mm}
        \end{pmatrix} \cdotp \begin{pmatrix}
        \sqrt{p_{1}^{0}}\\
        \sqrt{p_{2}^{0}}\\
        \dotsc \\
        \sqrt{p_{m}^{0}}
        \end{pmatrix} =\begin{pmatrix}
        1\\
        0\\
        \dotsc \\
        0
        \end{pmatrix} .
    \end{equation*}
    Таким образом, 
    \begin{equation*}
    V\left( I_{m} -zz^{T}\right) V^{T} =I_{m} -\begin{pmatrix}
    1 & 0 & \dotsc  & 0\\
    0 & 0 & \dotsc  & 0\\
    \dotsc  & \dotsc  & \dotsc  & \dotsc \\
    0 & 0 & \dotsc  & 0
    \end{pmatrix} =\begin{pmatrix}
    0 & 0 & \dotsc  & 0\\
    0 & 1 & \dotsc  & 0\\
    \dotsc  & \dotsc  & \dotsc  & \dotsc \\
    0 & 0 & \dotsc  & 1
    \end{pmatrix} =:\tilde{I}_{m} .
    \end{equation*}
    По теореме о наследовании сходимости
    \begin{equation*}
        \Vert V\xi _{n}\Vert ^{2}\xrightarrow{d}\left\Vert \mathcal{N}\left( 0,\ \tilde{I}_{m}\right)\right\Vert ^{2} =\chi _{m-1}^{2} .
    \end{equation*}
    Но в силу ортогональности матрицы $\displaystyle V$ выполнено
    
    
    \begin{equation*}
        \Vert V\xi _{n}\Vert ^{2} =\Vert \xi _{n}\Vert ^{2} =\left\Vert \left(\sqrt{B}\right)^{-1}\sqrt{n}\left(\dfrac{\overline{\nu }}{n} -\overline{p}_{0}\right)\right\Vert ^{2} =\sum _{i=1}^{m}\dfrac{\left( \nu _{i} -np_{i}^{0}\right)^{2}}{np_{i}^{0}} =\hat{\chi }_{n}^{2} .
    \end{equation*}
\end{proof}